\chapter{Checker Command-Line Interface}
\label{sec:checker-interface}

The proof checker is an external tool, separate from the solver, that verifies proofs written in the \vnnlib{} proof format. 

\section{The \texttt{check} Command}

A compliant checker should support the \texttt{check} command, with the following general usage pattern:
\begin{lstlisting}[style=bash]
<checker> check 
	--query <filepath>
	--proof <filepath> 
	[--network name=<filepath>]
	[ --timeout <seconds>]
\end{lstlisting}
This command has two arguments:



\clOption{-{}-query}{The path of the \vnnlib{} file that defines the problem claimed to be unsatisfiable }
{-{}-query /path/to/query.vnnlib}
\clOption{-{}-proof}{The path of the proof file to be checked}{-{}-proof /path/to/proof.alethe}
Additionally, the proof checker should support the following additional options:
\clOption{-{}-network}{This option maps the name of one of the networks declared in the provided \vnnlib{} query file to its implementation as an ONNX model file. When the query file contains multiple \texttt{declare-network} declarations this option should be used multiple times -- once for each network declaration not marked with an \texttt{equal-to} declaration (see Section~\ref{sec:multiple-networks}).}{-{}-network classifier=/a/path/to/a/model.onnx}

\clOption{-{}-timeout}{Maximum time to spend on processing on the proof in an integer number of seconds.}{-{}-timeout 10}

\subsection{Outputs of the \texttt{check} Command}
 The output of the \texttt{check} should be reported on \texttt{stdout} and should be a single line consisting of one of the following options:
\begin{itemize}
	\item \texttt{timed-out} - The allocated time elapsed before the checking procedure terminated.
	\item \texttt{valid} - The checker terminated and certified the proof witnesses the unsatisfiability of the query.
	\item \texttt{invalid} - The checker terminated with at least one assumption that could not be imported from the initial network or \vnnlib{} query, or at least one proof step whose checking has failed. The checker may add information for better understanding which proofs steps did not pass the checking and the reasons for failure.
	\item \texttt{holey} (optional) - The checker terminated and certified all assumptions and non-holey proof steps. Therefore, the proof witnesses the unsatisfiability of the query if and only if the \texttt{hole} steps are proven.
\end{itemize}
In particular, a proof checker is required to be complete and thus, unlike a solver, cannot output an \texttt{unknown} result.

\subsection{Error Conditions of the \texttt{check} Command}
The \texttt{check} command \emph{must} error when encountering any of the solver error cases as described in Section~\ref{sec:verify-command-errors}. In addition, if the checker is incapable of providing a sound analysis of the proof over the declared types in the \vnnlib{} query, it should require an explicit consent from the user, by re-writing the query to use the \inlinevnn{real} type (see Section~\ref{sec:real-queries} for details). Error messages must be printed on \texttt{stderr}.





